\documentclass[a5paper,10pt,idxtotoc]{labbook}
% A5, small pocket-notebook trim size for a bench book that travels with you
% between instruments, as opposed to a full-page desk notebook.

\usepackage[utf8]{inputenc}
\usepackage[T1]{fontenc}
\usepackage{booktabs}
\usepackage{amsmath}
\usepackage[margin=1.6cm]{geometry}

\newexperiment{calib}{Oscilloscope probe calibration}
\newexperiment{filter}{Low-pass filter characterisation}
\newexperiment{amp}{Op-amp gain stage}

\title{Bench Notebook --- Analogue Filter Board}
\author{R.\ Voss}
\date{Notebook 3}

\begin{document}

\frontmatter
\title{Bench Notebook --- Analogue Filter Board}
\author{R.\ Voss}
\date{Notebook 3}
\maketitle
\printindex
\tableofcontents

\mainmatter

\labday{9 June}

\experiment{calib}
Compensated both scope probes against the built-in 1\,kHz square wave before
starting: probe~1 needed a small trim (slightly overcompensated out of the
box), probe~2 was fine as received.

\experiment{filter}
Built the first-pass RC low-pass on the breadboard: R = 1.6\,k$\Omega$,
C = 100\,nF, giving a calculated corner of

\[
  f_c = \frac{1}{2\pi RC} \approx 995\text{ Hz}.
\]

Swept the input from 100\,Hz to 10\,kHz at 1\,V\textsubscript{pp} and logged
the output amplitude:

\begin{table}[h]
  \centering
  \begin{tabular}{rr}
    \toprule
    Freq. (Hz) & Output (V\textsubscript{pp}) \\
    \midrule
    100  & 0.99 \\
    500  & 0.94 \\
    1000 & 0.70 \\
    2000 & 0.44 \\
    5000 & 0.19 \\
    10000 & 0.10 \\
    \bottomrule
  \end{tabular}
  \caption{Filter sweep, first breadboard pass}
\end{table}

Measured corner (--3\,dB point, $0.707\,V_{pp}$) lands close to 1\,kHz,
consistent with the calculated value.

\labday{10 June}

\experiment{amp}
Added the op-amp gain stage after the filter. Target gain 10 V/V using
R\textsubscript{f} = 10\,k$\Omega$, R\textsubscript{g} = 1\,k$\Omega$ in a
non-inverting configuration.

Measured DC offset at the output with no input: 4\,mV, well within spec for
this part. Gain measured at 1\,kHz: 9.7 V/V, close enough to the 10 V/V
target that the 1\% resistor tolerance accounts for the gap.

\experiment{filter}
Re-ran the sweep with the amplifier stage in place to check for interaction
between the two stages. Corner frequency unchanged; the amplifier's input
impedance is high enough not to load the filter.

\textit{Follow-up:} move from breadboard to a soldered prototype before the
next set of measurements --- breadboard parasitic capacitance may be
softening the roll-off above 5\,kHz.

\labday{12 June}

\experiment{amp}
Soldered prototype on stripboard. Re-measured gain: 9.8 V/V, DC offset
3\,mV. Noise floor visibly lower than the breadboard version on the scope.

\experiment{calib}
Probe~1 drifted slightly since the 9th; re-trimmed compensation before
today's measurements.

\end{document}
